\SourceLeaf{181}are always greater than the assignable difference between the hypotenuse and one
leg of \textls{the particular} right triangle in which each in turn is the
other leg. Nevertheless, their sum plainly never exhausts the
hypotenuse $AC$, however many of them are taken.

We would almost have to regret dwelling so long on matters so familiar and simple
if this work, whose author, incidentally, does seem sincerely concerned with
the truth, had not drawn more than the usual attention by the manner in which it
had been announced in the public press even before its
appearance. We shall only add here that
the author subsequently attempts, by a very similar and therefore equally futile method, to prove
that the angle $MST$ cannot be obtuse. Here, however, there is an
essential difference, for the impossibility of this case can in fact be proved with
complete rigor, although this is not the place to develop the proof
further.

\Needspace{11\baselineskip}
\begin{center}\rule{3cm}{0.4pt}\end{center}

\begin{center}
\textsc{Remark.}
\end{center}

This review, already reprinted in Volume IV, pages 364 to 368, is reproduced here
for completeness. Two figures, which appear neither in the
original nor in the first reprint, are intended to make the text
easier to understand.

\begin{flushright}\textsc{Stäckel.}\end{flushright}

\begin{center}\rule{3cm}{0.4pt}\end{center}
